\documentclass[10.5pt,letterpaper]{article}

\usepackage[letterpaper,margin=0.82in,headheight=15pt]{geometry}
\usepackage[T1]{fontenc}
\usepackage{lmodern}
\usepackage[scaled=0.92]{helvet}
\usepackage{inconsolata}
\usepackage{microtype}
\usepackage{xcolor}
\usepackage{graphicx}
\usepackage{booktabs}
\usepackage{tabularx}
\usepackage{longtable}
\usepackage{array}
\usepackage{enumitem}
\usepackage{amsmath}
\usepackage{amssymb}
\usepackage{mathtools}
\usepackage{listings}
\usepackage[most]{tcolorbox}
\usepackage{titlesec}
\usepackage{fancyhdr}
\usepackage{hyperref}
\usepackage{bookmark}
\usepackage{needspace}
\usepackage{ragged2e}
\usepackage{lastpage}

\definecolor{navy}{HTML}{19344A}
\definecolor{blue}{HTML}{246FA8}
\definecolor{deepblue}{HTML}{164E73}
\definecolor{gold}{HTML}{A36D00}
\definecolor{ink}{HTML}{20252B}
\definecolor{muted}{HTML}{626B73}
\definecolor{paleblue}{HTML}{EAF2F8}
\definecolor{palegold}{HTML}{FBF3DE}
\definecolor{palegray}{HTML}{F4F6F7}
\definecolor{lineblue}{HTML}{AFC9DD}
\definecolor{good}{HTML}{2B6B45}
\definecolor{warn}{HTML}{8C5A00}

\hypersetup{
  colorlinks=true,
  linkcolor=deepblue,
  urlcolor=blue,
  pdftitle={Why This AlphaZero Loop Worked},
  pdfauthor={Breakthrough AlphaZero project},
  pdfsubject={Technical post-mortem and implementation guide}
}

\pagestyle{fancy}
\fancyhf{}
\fancyhead[L]{\sffamily\footnotesize\color{muted}ALPHAZERO POST-MORTEM}
\fancyhead[R]{\sffamily\footnotesize\color{muted}BREAKTHROUGH PROJECT}
\fancyfoot[L]{\sffamily\footnotesize\color{muted}Technical implementation guide}
\fancyfoot[R]{\sffamily\footnotesize\color{muted}\thepage\ / \pageref{LastPage}}
\renewcommand{\headrulewidth}{0.35pt}
\renewcommand{\footrulewidth}{0pt}

\titleformat{\section}
  {\needspace{4\baselineskip}\sffamily\Large\bfseries\color{blue}}
  {\thesection}{0.65em}{}
\titleformat{\subsection}
  {\needspace{3\baselineskip}\sffamily\large\bfseries\color{deepblue}}
  {\thesubsection}{0.6em}{}
\titleformat{\subsubsection}
  {\needspace{2.5\baselineskip}\sffamily\normalsize\bfseries\color{navy}}
  {\thesubsubsection}{0.55em}{}
\titlespacing*{\section}{0pt}{1.6em}{0.65em}
\titlespacing*{\subsection}{0pt}{1.2em}{0.45em}
\titlespacing*{\subsubsection}{0pt}{0.9em}{0.3em}

\setlength{\parindent}{0pt}
\setlength{\parskip}{0.55em}
\setlength{\tabcolsep}{5pt}
\renewcommand{\arraystretch}{1.22}
\setlist[itemize]{leftmargin=1.35em,itemsep=0.3em,topsep=0.3em}
\setlist[enumerate]{leftmargin=1.6em,itemsep=0.45em,topsep=0.35em}

\newcolumntype{Y}{>{\RaggedRight\arraybackslash}X}
\newcolumntype{L}[1]{>{\RaggedRight\arraybackslash}p{#1}}

\newtcolorbox{keybox}[1][]{
  enhanced, colback=paleblue, colframe=blue, boxrule=0pt,
  borderline west={2.4pt}{0pt}{blue}, arc=0pt,
  left=8pt,right=7pt,top=6pt,bottom=6pt,
  before skip=8pt,after skip=9pt, #1
}
\newtcolorbox{warningbox}[1][]{
  enhanced, colback=palegold, colframe=gold, boxrule=0pt,
  borderline west={2.4pt}{0pt}{gold}, arc=0pt,
  left=8pt,right=7pt,top=6pt,bottom=6pt,
  before skip=8pt,after skip=9pt, #1
}
\newtcolorbox{plainbox}[1][]{
  enhanced, colback=palegray, colframe=lineblue, boxrule=0.5pt,
  arc=1pt, left=8pt,right=8pt,top=6pt,bottom=6pt,
  before skip=8pt,after skip=9pt, #1
}

\newcommand{\boxtitle}[1]{\sffamily\bfseries\color{deepblue}#1.\ }
\newcommand{\code}[1]{\texttt{\small #1}}

\lstdefinestyle{azcode}{
  basicstyle=\ttfamily\small\color{navy},
  backgroundcolor=\color{palegray},
  frame=single,
  rulecolor=\color{lineblue},
  framesep=6pt,
  xleftmargin=8pt,
  xrightmargin=4pt,
  aboveskip=7pt,
  belowskip=8pt,
  showstringspaces=false,
  columns=fullflexible,
  keepspaces=true
}

\begin{document}

\begin{titlepage}
\thispagestyle{empty}
\vspace*{0.65in}
{\sffamily\bfseries\color{gold}\large TECHNICAL POST-MORTEM AND IMPLEMENTATION GUIDE}
\vspace{0.55in}

{\sffamily\bfseries\color{navy}\fontsize{31}{35}\selectfont
Why This AlphaZero Loop Worked\par}
\vspace{0.25in}

{\sffamily\color{deepblue}\fontsize{17}{22}\selectfont
Lessons from repeated failures, a clean-room rebuild,\\
and a successful 8x8 Breakthrough run\par}
\vspace{0.28in}

{\large\itshape\color{muted}
For students implementing AlphaZero-style systems\\
26 August 2026\par}

\vfill
\begin{keybox}
\boxtitle{Project-declared success}
The accepted iteration-26 network beat the already-trained phase-6 starting
model 39--1, then beat iteration 15 by 36--4, and finally beat the depth-4
alpha-beta baseline 49--1 in a 50-game paired-opening arena. The 50-hour
training program was still running when this post-mortem was written.
\end{keybox}
\vspace{0.2in}

\begin{center}
\large\itshape
The lesson is not that one hyperparameter was magic. The system improved when
its conventions, tests, data, search, optimization, exploration, and evaluation
finally formed one coherent feedback loop.
\end{center}

\vfill
{\sffamily\footnotesize\color{muted}
Breakthrough AlphaZero project \hfill Post-mortem and student guide}
\end{titlepage}

\setcounter{tocdepth}{1}
\tableofcontents
\clearpage

\section*{How to read this post-mortem}
\addcontentsline{toc}{section}{How to read this post-mortem}

Begin with the failure mode, not the final architecture. Suppose a Player-2
node mistakenly treats a positive Player-1 value as desirable. Search then
spends visits on the wrong move. Those visits become the policy target. The
network learns to reproduce the mistake, and the next round of self-play makes
the same data error more confidently. Nothing crashes. The losses may even
fall. The circular pipeline has simply learned to agree with itself.

This is the central difficulty of AlphaZero. The policy and value equations are
compact; the hard part is keeping their meanings unchanged as positions pass
through rules, coordinate transforms, search, replay, training, checkpoints,
and evaluation. This chapter follows those meanings all the way around the
loop. By the end, a reader should be able to:

\begin{itemize}
  \item derive selection and backup from a stated value convention;
  \item trace one self-play position into its policy and value targets;
  \item design tests for rules, actions, symmetry, masking, and serialization;
  \item distinguish a search failure from a learning or throughput failure;
  \item design a paired arena that measures strength rather than randomness;
  \item adapt the method without importing Breakthrough-specific assumptions.
\end{itemize}

Reinforcement-learning success also invites retrospective storytelling: after
the score improves, every design choice can look inevitable. This guide
therefore separates three kinds of statement.

\begin{center}
\small
\begin{tabularx}{\textwidth}{L{0.19\textwidth} Y Y}
\toprule
\sffamily\bfseries Evidence class &
\sffamily\bfseries Meaning &
\sffamily\bfseries Example in this project \\
\midrule
Invariant &
A property established by unit, differential, or metamorphic tests. &
Action round trips, perspective conversion, sign-free backup, and exact
make/unmake restoration. \\
Measured result &
A comparison from a completed run with recorded settings. &
The 1,000-game versus 10,000-game pretraining match and paired checkpoint
arenas. \\
Causal hypothesis &
A plausible explanation consistent with the evidence, but not isolated by a
full ablation. &
Why the combined serious-training recipe escaped the feedback failures of
earlier attempts. \\
\bottomrule
\end{tabularx}
\end{center}

\begin{warningbox}
\boxtitle{Causality warning}
The successful 8x8 run changed several compute-efficiency and training choices
together. Its success proves that the complete recipe works in this setting.
It does not prove that every ingredient is necessary or that each contributed
independently. Where the evidence does not support a causal claim, this guide
says so.
\end{warningbox}

\begin{plainbox}
\boxtitle{The guide in one sentence}
First make every boundary testable; then give the learner enough diverse,
search-improved data; finally measure strength with paired games rather than
with training loss.
\end{plainbox}

\subsection*{Three useful routes through the chapter}

\begin{center}
\small
\begin{tabularx}{\textwidth}{L{0.22\textwidth} Y}
\toprule
\sffamily\bfseries If you are\ldots &
\sffamily\bfseries Read in this order \\
\midrule
Implementing the system &
Sections 3--6, then the staged protocol in Section 9. Treat every displayed
invariant as a test requirement. \\
Diagnosing a run that does not improve &
The worked failure in Section 3, the 0--32 case study in Section 6, and the
symptom-to-test table in Section 7. \\
Adapting the project to another environment &
The value, state, and symmetry contracts in Section 4, then all of Section 8.
Return to the Breakthrough recipe only for examples, not default numbers. \\
\bottomrule
\end{tabularx}
\end{center}

\section{What succeeded}

The project succeeded in the operational sense that matters for an AlphaZero
assignment: repeated
\[
  \textsc{Play} \longrightarrow \textsc{Search} \longrightarrow
  \textsc{Train}
\]
updates produced agents demonstrably stronger than their trained predecessors
and stronger than a fixed classical baseline. That is much stronger evidence
than a falling loss, a plausible-looking policy, or a single entertaining game.

\begin{center}
\small
\begin{tabularx}{\textwidth}{Y L{0.12\textwidth} Y}
\toprule
\sffamily\bfseries Checkpoint comparison &
\sffamily\bfseries Score &
\sffamily\bfseries Protocol \\
\midrule
Phase-6 iteration 5 vs.\ raw pretrained initializer &
52--8 &
60 paired 0.1-second games. \\
\addlinespace[2pt]
Serious iteration 15 vs.\ phase-6 iteration 5 &
39--1 &
20 six-ply openings, colors reversed, fixed 64-simulation search. \\
\addlinespace[2pt]
Serious iteration 26 vs.\ serious iteration 15 &
36--4 &
The same scheduled strength-check protocol. \\
\addlinespace[2pt]
Serious iteration 26 vs.\ depth-4 alpha-beta &
49--1 &
25 six-ply openings, colors reversed, 0.1 seconds per move. \\
\bottomrule
\end{tabularx}
\end{center}

The final alpha-beta match was balanced by color: iteration 26 scored 24--1 as
Player 1 and 25--0 as Player 2. All 50 games completed. Its Wilson interval was
89.5\%--99.6\%. The match ran on an L40S, whereas the earlier 2--58 phase-6
result used RTX 2080-class hardware. The exact magnitude of improvement is
therefore not hardware-controlled. The conclusion that iteration 26 beat the
baseline under the measured protocol is nevertheless unambiguous.

\subsection*{What success does not establish}
\begin{itemize}
  \item It does not establish optimal play, or even the strongest feasible
        Breakthrough agent.
  \item It does not prove that the full 50-hour budget was necessary; the agent
        was already strong by iteration 26.
  \item It does not isolate the individual contribution of playout-cap
        randomization, replay size, root noise, search growth, or learning rate.
  \item It does not make validation loss a strength metric. The decisive
        evidence is the checkpoint ladder and the external baseline match.
  \item It does not imply that the same numeric settings transfer unchanged to
        games with different branching factors, horizons, rewards, or
        symmetries.
\end{itemize}

\section{Why earlier attempts failed}

The earlier projects did not fail because they lacked sophistication. They
failed because individually reasonable components disagreed at their
boundaries. A process-local augmentation counter reset. An inference ensemble
required by training was omitted. Rare teacher samples dominated batches.
Optimizer and checkpoint lifecycles became unclear. Some arenas compared
identical model hashes or counted time forfeits as strength.

AlphaZero is unusually unforgiving of such defects because each model creates
the next model's data. A small systematic error is not merely observed; it is
fed back into the system.

\small
\begin{longtable}{L{0.19\textwidth} L{0.38\textwidth} L{0.35\textwidth}}
\toprule
\sffamily\bfseries Earlier pattern &
\sffamily\bfseries Observed failure &
\sffamily\bfseries Successful-project response \\
\midrule
\endfirsthead
\toprule
\sffamily\bfseries Earlier pattern &
\sffamily\bfseries Observed failure &
\sffamily\bfseries Successful-project response \\
\midrule
\endhead
Complex absolute-player encoding plus four transformations &
A prior 8x8 continuous run lost augmentation state between short learner
processes and omitted the required inference ensemble. The actor developed
asymmetric blind spots and material gifts. &
Use two mover-relative planes. Convert value perspective in one wrapper.
Player swap becomes a test, not training augmentation; random left-right
reflection is enough. \\
\addlinespace
Tiny or badly weighted learning data &
One historical diagnosis gave rare teacher positions about 71 times the
per-position weight of neural positions. A nominal batch of 256 had effective
size near 18, and every trained epoch worsened validation. &
Use a recent FIFO of homogeneous neural self-play for reinforcement learning.
Size training by presentations per fresh position, not opaque source weights. \\
\addlinespace
Optimizer and actor lifecycle spread across jobs &
Restarted learner processes and rollback machinery made it easy to lose Adam
state or compare one checkpoint under two labels. &
Save Keras optimizer state with every checkpoint. The newest trained network
always generates the next tranche; reference checkpoints are diagnostic only. \\
\addlinespace
Repeated live seeds &
Supposedly fresh runs reproduced the same first games. Deterministic evaluation
risked measuring one repeated trajectory. &
Do not fix seeds in production self-play, replay sampling, or arenas. Generate
distinct random opening prefixes, then reverse colors within each pair. \\
\addlinespace
Too little search mistaken for a batching bug &
At 8/2 simulations, sequential and batched self-play both produced 0--32
Player-1 results. &
Separate algorithm quality from throughput. A 64-simulation diagnostic restored
a healthy 19--13 split; the serious schedule began at 128/16. \\
\addlinespace
Metrics and gates accumulated faster than understanding &
Binary tactical gates and actor-versus-candidate machinery sometimes rejected
useful models or evaluated identical hashes. &
Keep direct invariant tests, a continuous tactical score, compact metrics, and
standalone paired arenas. \\
\bottomrule
\end{longtable}
\normalsize

\begin{keybox}
\boxtitle{Most important comparison}
The historical failures came from sophistication crossing untested boundaries:
process-local state, replay weights, time-control grace, checkpoint identity,
symmetry assumptions, and actor lifecycle. The successful rebuild deliberately
reduced the number of boundaries and made the remaining ones observable.
\end{keybox}

\section{Follow one bug around the loop}

The most dangerous bugs in AlphaZero are semantic: every array has the right
shape, every probability sums to one, and the program trains. To see why they
matter, consider a search tree whose $Q$ values are always measured from
Player 1's viewpoint.

At a Player-2 parent, two moves have the following statistics. The exploration
bonus has already been calculated from the prior and visit counts.

\begin{center}
\small
\begin{tabularx}{0.78\textwidth}{L{0.15\textwidth} Y Y Y}
\toprule
\sffamily\bfseries Move &
\sffamily\bfseries Absolute $Q$ &
\sffamily\bfseries Exploration $U$ &
\sffamily\bfseries Meaning \\
\midrule
$A$ & $+0.45$ & $0.12$ & Strongly favorable to Player 1. \\
$B$ & $-0.10$ & $0.12$ & Slightly favorable to Player 2. \\
\bottomrule
\end{tabularx}
\end{center}

If selection blindly maximizes $Q+U$, it gives $A$ a score of $0.57$ and $B$
a score of $0.02$. Player 2 selects $A$, the move that is better for the
opponent. With an absolute Player-1 value, the parent must maximize
$p(s)Q+U$. Since $p(s)=-1$ here, the scores are
\[
  \operatorname{score}(A)=-0.45+0.12=-0.33,
  \qquad
  \operatorname{score}(B)=+0.10+0.12=+0.22.
\]
The corrected search selects $B$.

Now follow the wrong calculation around one training cycle:
\begin{enumerate}
  \item PUCT visits $A$ too often at Player-2 nodes.
  \item The normalized visit counts assign a high policy target to $A$.
  \item Self-play chooses $A$ often and produces states from the wrong branch.
  \item Gradient descent faithfully teaches the policy to prefer $A$ in similar
        positions.
  \item The next search begins with a larger prior on $A$, amplifying the same
        mistake.
\end{enumerate}

This explains why ``the loss went down'' is weak evidence. The network can
learn corrupted targets perfectly. The antidote is not a different optimizer;
it is a hand-computed search test that includes both players.

\subsection{The same example through the value boundary}

Suppose the network evaluates a leaf where Player 2 is to move and returns
$v_{\mathrm{rel}}=+0.70$. The prediction is favorable to the mover, so in the
tree's absolute coordinate system it means
\[
  v_{\mathrm{abs}}=(-1)(+0.70)=-0.70.
\]
Every edge on the search path receives the same absolute result $-0.70$;
backup does not change its sign. Only selection interprets that result for the
parent whose choice is being made. At a Player-1 parent it contributes
$+1\cdot(-0.70)$; at a Player-2 parent it contributes
$-1\cdot(-0.70)=+0.70$.

\begin{keybox}
\boxtitle{The invariant to put beside the code}
$W$ and $Q$ answer one fixed question: ``How good is this for Player 1?''
The network wrapper converts viewpoints once. Backup averages absolute values
without sign changes. Selection multiplies exploitation by the \emph{parent's}
player. If any module answers a different question, make the conversion visible
at that boundary.
\end{keybox}

\section{The contracts that made the system coherent}

\subsection{One public value convention}

Outside the neural network, every value is an \emph{absolute Player-1 value}:
$+1$ favors Player 1 and $-1$ favors Player 2. The network sees the mover as
its own side, so its raw value is mover-relative. A single neural-boundary
wrapper performs the conversion:
\[
  v_{\mathrm{abs}}(s) = p(s)\,v_{\mathrm{rel}}(s),
  \qquad
  z_{\mathrm{rel}}(s) = p(s)\,z_{\mathrm{abs}},
\]
where $p(s)=+1$ when Player 1 is to move and $p(s)=-1$ when Player 2 is to move.

The tree stores absolute values. Backup therefore never flips a sign. Selection
is the only turn-dependent operation:
\[
  a^* = \arg\max_a
  \left[
    p(s)\,Q(s,a) +
    c_{\mathrm{puct}}\,P(s,a)
    \frac{\sqrt{\sum_b N(s,b)}}{1+N(s,a)}
  \right].
\]
Player 1 selects larger absolute $Q$; Player 2 selects smaller absolute $Q$.
Tests exercise both players on hand-computed trees.

\needspace{13\baselineskip}
In code, the two places where the convention appears can remain this direct:

\begin{lstlisting}[style=azcode]
exploration = cpuct * child.prior * sqrt(parent.visits) \
              / (1 + child.visits)
score = parent_player * child.q + exploration

while node is not None:
    node.visits += 1
    node.q += (absolute_value - node.q) / node.visits
    node = node.parent
\end{lstlisting}

There is no sign change in the backup because \code{absolute\_value} keeps the
same meaning at every depth. If the alternative mover-relative convention is
used, derive its backup separately rather than mixing the two snippets.

An unvisited child begins with its parent's current $Q$. This is the project's
first-play estimate: an unseen action is not treated as a mysterious zero when
the parent already has evidence. When a new leaf is evaluated, back it up
before expanding it. Its first incremental-mean update then sets its $Q$ exactly
to the leaf value, and the children it creates inherit that estimate. The
ordering is a small code detail with a precise statistical meaning, so it
deserves a direct unit test.

\begin{plainbox}
\boxtitle{Tradeoff}
Absolute tree values are not inherently better than mover-relative tree values.
Mover-relative values can be elegant, but then every edge crossing changes
perspective and backup usually flips a sign. Absolute values made the teaching
implementation easier to audit because one wrapper and one exploit term contain
all perspective logic. The portable rule is not to always use absolute values.
It is to choose one convention, write it down, and test every boundary against
it.
\end{plainbox}

\subsection{A trusted ladder below self-play}

The project did not ask a neural reinforcement-learning loop to diagnose its
own foundations. Each layer had a cheaper oracle:
\begin{itemize}
  \item \textbf{Rules:} an independently written legal-move generator and exact
        make/unmake restoration.
  \item \textbf{Terminal logic:} goal-row, no-reply, capture, and the rule that
        a terminal move does not switch player.
  \item \textbf{Action encoding:} encode/decode round trips for every legal move
        in random games.
  \item \textbf{Symmetry:} metamorphic tests for legal moves, policy indices,
        outcomes, and reflected positions.
  \item \textbf{Search:} uniform-prior or rollout search that finds immediate
        wins for both players before a neural network is introduced.
  \item \textbf{Network boundary:} solver-labelled supervised positions that
        establish that policy masking, value perspective, fitting, save/load,
        and inference all agree.
\end{itemize}

This ladder matters because a self-play loop that did not improve is an
extremely low-information symptom. A reflected position that changes its
terminal outcome, or a Player-2 node that prefers the larger absolute $Q$,
points almost directly to a faulty line of code.

\subsection{Canonical state and symmetry}

The neural input contains two planes: pieces belonging to the mover and pieces
belonging to the opponent. The board is oriented so the mover always advances
in the same neural direction. This removes an irrelevant color distinction
without discarding game information.

For Breakthrough, player swap is already represented by mover-relative
canonicalization, so applying it again during training adds no new example.
Left-right reflection is the genuine geometric symmetry. The replay buffer
stores the original position; a random reflection is applied after sampling,
with the policy target transformed by the same action permutation.

For an $n\times n$ board, reflection sends a move tuple
\[
  (r_1,c_1,r_2,c_2)
  \longmapsto
  (r_1,n-1-c_1,r_2,n-1-c_2).
\]
On an 8x8 board, $(2,1,3,2)$ therefore becomes $(2,6,3,5)$. In the three
policy channels attached to every source square, forward-left and
forward-right exchange places while forward stays fixed. Thus a one-hot target
for the original move must become a one-hot target for the reflected move; it
must not remain at the same flat policy index.

This gives a useful metamorphic test. Start from a legal position $s$, reflect
it to $R(s)$, and reflect every legal action. For every action $a$ require
\[
  R(T(s,a)) = T(R(s),R(a)),
\]
where $T$ is the game transition. Then transform a policy vector, invert the
transformation, and require exact recovery of the original vector.

\begin{warningbox}
\boxtitle{Common symmetry bug}
Transforming the board but not the policy indices silently assigns probability
to the wrong actions. Test a one-hot policy through every transformation and its
inverse. Also test legal-action masks and terminal rewards. A visual inspection
of transformed boards is not enough.
\end{warningbox}

\subsection{Exact state restoration}

MCTS applies and retracts many moves. A one-bit mutation leak can corrupt
distant branches while leaving the root position apparently plausible. For
random legal positions and every legal action, the project checks:
\[
  \texttt{snapshot(state)}
  =
  \texttt{snapshot(unmake(make(state, action)))}.
\]
The snapshot includes the board, player to move, terminal fields, and any cached
legal-action information. For a new environment, hash the complete state before
and after every debug search until this contract is beyond doubt.

\section{Data and optimization choices that mattered}

\subsection{The 5x5 ladder was an instrument, not a toy}

The smaller board shortened games, reduced the branching factor, and made exact
or strong tactical labels affordable. It exposed value-perspective errors,
BatchNorm behavior, memorization, label imbalance, and weaknesses in the first
tactical suite before expensive 8x8 jobs were submitted.

The project did \emph{not} stretch a padded 5x5 network into an 8x8 network and
call that evidence about native 8x8 strength. Native board sizes received
native models and native data. What transferred was the tested process.

\begin{plainbox}
\boxtitle{Generalization}
Choose the smallest faithful instance of the target environment: a smaller
board, shorter horizon, reduced action set, or deterministic scenario. It must
preserve the contracts you are testing. If the reduced problem changes the
meaning of an action, observation, or reward, it is a different environment,
not a debugging ladder.
\end{plainbox}

\subsection{Pretraining gave self-play a competent starting distribution}

The 8x8 initializer was trained on 10,000 rollout-PUCT games containing 421,561
stored positions. A 1,000-game subset was not nearly as good: it lost 24--76 to
the 10,000-game model. The large corpus therefore bought measurable strength,
not merely a prettier loss curve.

The stored corpus contained real positions. Symmetry was sampled during
training rather than materialized as duplicate records. This keeps storage
honest and allows a fresh random transformation on every presentation.

Pretraining is not a requirement of AlphaZero in principle. In a modest compute
budget it is an engineering choice: it avoids spending the first expensive
neural self-play hours rediscovering elementary legality and tactics. For
another environment, a teacher may be a rollout search, alpha-beta, heuristic
controller, solver, demonstration set, or a mixture. The target is not
imitation forever; it is a policy/value pair good enough for search to produce
useful improvements.

\subsection{Architecture followed evidence}

The final network was intentionally modest: 48 filters, three residual blocks,
and a direct three-channel $1\times1$ convolutional policy head. Batch
normalization was removed after supervised diagnostics showed that its moving
statistics and small, correlated batches complicated the learning signal.

This is a project-specific result, not a universal ban on BatchNorm. Larger,
well-shuffled batches or different normalization schemes may work well. The
portable lesson is to ask the network to pass a controlled supervised task
before blaming self-play. If a small model cannot fit and generalize on exact
labels, more actors will only generate expensive ambiguity.

\subsection{Replay exposure was explicit}

The serious loop used a FIFO replay of roughly 25,000 recent neural positions,
batch size 128, and about two training presentations per fresh position. Adam
state was retained, with learning rate $2.5\times10^{-4}$.

The useful quantity is not epochs in isolation. It is how often each fresh
position is expected to influence an update:
\[
  \text{reuse ratio}
  =
  \frac{\text{training examples drawn in an iteration}}
       {\text{new positions added in that iteration}}.
\]
Too little reuse underfits expensive search targets. Too much reuse memorizes a
small, highly correlated window and can erase older competencies. The right
window also depends on policy drift: recent data are on-policy but narrow;
older data improve diversity but were generated by weaker policies.

\begin{warningbox}
\boxtitle{Do not fix divergence by freezing learning}
Lowering the learning rate to nearly zero can make losses and weights look
stable while stopping improvement. First diagnose perspective, labels,
effective sample size, replay composition, normalization, optimizer state, and
search quality. Stability without learning is not success.
\end{warningbox}

\subsection{The loss was conventional; the targets were the hard part}

For canonical state $s$, search policy $\pi$, and mover-relative outcome $z$,
the learner minimized
\[
  \mathcal{L}(\theta)
  =
  (z-v_\theta(s))^2
  -
  \pi^\mathsf{T}\log p_\theta(\cdot\mid s)
  +
  \lambda\lVert\theta\rVert_2^2.
\]
Illegal logits are masked before a policy is used. Search visit counts are
normalized over legal root actions to create $\pi$. The value target is the
eventual game result, converted exactly once at the neural boundary.

Most catastrophic loss problems were not caused by this equation. They came
from feeding the equation targets expressed in the wrong coordinate system,
perspective, action indexing, or sample distribution.

\subsection{Validation answers a narrower question than an arena}

For the fixed pretraining corpus, complete games were assigned to either the
training split or the validation split before positions were extracted and
augmented. Splitting individual positions would leak near-duplicates from the
same trajectory across both sets. Applying reflection before the split could
leak exact symmetric copies.

A validation set asks whether the model generalizes to held-out examples from
the same data-producing process. It can expose memorization, a representation
mistake, or a normalization/serialization discrepancy. It cannot, by itself,
answer whether the policy wins more games. During self-play the actor and its
state distribution are changing, so even the meaning of a stationary
validation curve becomes limited.

\begin{center}
\small
\begin{tabularx}{\textwidth}{Y Y}
\toprule
\sffamily\bfseries Question & \sffamily\bfseries Appropriate evidence \\
\midrule
Can the network learn the teacher's labels beyond its training games? &
Whole-game held-out policy and value losses. \\
Are state, action, and value targets internally correct? &
Exact labels, round trips, metamorphic tests, and hand calculations. \\
Was the larger corpus worth its generation cost? &
A data-size curve or a match between models trained on nested subsets. \\
Did the latest checkpoint become a stronger player? &
A controlled arena using search, paired openings, and both colors. \\
\bottomrule
\end{tabularx}
\end{center}

\section{A serious but simple AlphaZero loop}

\subsection{The loop}

The production design remained conceptually small:

\begin{lstlisting}[style=azcode]
while time_is_left:
    games = play_parallel_games(current_model)
    replay.add(games)

    train_current_model(replay)
    save_model_and_optimizer()

    if strength_check_is_due:
        play_paired_match(current_model, reference_model)
\end{lstlisting}

There is no candidate-acceptance gate in the learning path. The latest trained
model generates the next tranche. Scheduled matches report progress; they do
not roll the actor backward. This avoids a noisy gate freezing learning or
creating a two-model lifecycle that students cannot audit.

\subsection{One training record, end to end}

The clearest way to understand the loop is to trace one position. While a game
is in progress, keep the untransformed board, the player to move, and the root
visit counts. The winner is not known yet:

\begin{lstlisting}[style=azcode]
history.append((board.copy(), player, root_visit_counts))
\end{lstlisting}

When the game ends, attach the same absolute outcome to every history entry:

\begin{lstlisting}[style=azcode]
for board, player, counts in history:
    replay.add((board, player, counts, winner_for_player_1))
\end{lstlisting}

Only after a record is sampled for training do we optionally reflect both its
board and visit-count vector. The neural input is then oriented to the mover,
and the absolute game outcome is converted to the mover's viewpoint:

\begin{lstlisting}[style=azcode]
board, player, counts, z_absolute = replay.sample_one()
if random.choice([False, True]):
    board = reflect_board(board)
    counts = reflect_policy(counts)
x = canonical_input(board, player)
pi = counts / counts.sum()
z = player * z_absolute
train_on(x, pi, z)
\end{lstlisting}

Here \code{player} is $+1$ for Player 1 and $-1$ for Player 2. The policy
target comes from search, not from the move eventually sampled for play. The
value target comes from the completed game, not the leaf evaluation used
during search. These two distinctions are easy to lose in a larger actor.

\begin{plainbox}
\boxtitle{Audit a replay row by hand}
Choose one stored row. Reconstruct its legal actions, normalize its visit
counts, reflect it and invert the reflection, compute its mover-relative value
target, and feed it through the saved model. This single exercise crosses most
of the interfaces that otherwise fail silently.
\end{plainbox}

\subsection{Parallel search changed throughput, not semantics}

Actors advanced 32 games together and batched neural leaf evaluations on the
GPU. Within one game, PUCT selection, expansion, backup, and root targets kept
their ordinary meaning. Batching was treated as a scheduling optimization
rather than a new search algorithm.

This distinction paid off during diagnosis. A severe Player-1 collapse appeared
at an 8/2 simulation budget in both sequential and batched implementations.
Raising the budget to 64 simulations restored a 19--13 split. The defect was
insufficient search, not batching.

\begin{warningbox}
\boxtitle{Worked diagnosis: the 0--32 result}
\textbf{Observation:} the batched actor produced 32 Player-2 wins and no
Player-1 wins. \textbf{First hypothesis:} batching had changed search
semantics. \textbf{Discriminating test:} run the slow sequential actor with the
same tiny 8/2 full/fast budget. It also produced 0--32, so the failure followed
the budget rather than the scheduler. \textbf{Second test:} keep the checkpoint
fixed and raise search to 64 simulations. The split became 19--13.

This did not prove that batching was bug-free; the reference comparison tests
did that separately. It showed that this particular collapse was caused by
asking an asymmetric game to generate its own labels with almost no lookahead.
The experiment changed one causal variable at a time and used an outcome that
could distinguish the hypotheses.
\end{warningbox}

\begin{plainbox}
\boxtitle{Portable principle}
Build and test a slow reference search first. The optimized actor should agree
with it on legal moves, terminal values, visit accounting, and fixed-network
decisions. Benchmark speed only after semantic equivalence is established.
\end{plainbox}

\subsection{Playout-cap randomization spent search where labels mattered}

Not every self-play move was stored. A minority of positions received the full
search budget and became training targets; intervening moves used a smaller
budget to advance games cheaply. In the serious schedule, about 25\% of turns
were full-search targets.

This trades target density for target quality and game throughput. It works
when the cheap turns remain plausible enough to reach useful states and the
full-search positions cover the game. Log target ply, player, outcome, and game
phase. If full searches cluster accidentally, the replay buffer will inherit
that blind spot.

\subsection{Search capacity grew with the model}

The schedule increased full/fast simulations from roughly 128/16 to 192/24 and
then 256/32. Early networks do not justify the same search expense as later
ones, but extremely shallow search can create self-confirming nonsense.

A portable calibration procedure is:
\begin{enumerate}
  \item On a fixed checkpoint, sweep several search budgets.
  \item Measure throughput, color balance, tactical score, policy entropy, and
        head-to-head strength.
  \item Choose the smallest budget before the strength curve clearly flattens.
  \item Repeat after the network becomes materially stronger.
\end{enumerate}

\subsection{Exploration in self-play; controlled diversity in evaluation}

Self-play added Dirichlet noise to root priors and sampled from visit counts
during the first 12 plies. Later moves were selected more sharply. This creates
varied openings while allowing games to become decisive.

Evaluation removed root noise and stochastic move sampling. Diversity came
from independently generated six-ply opening prefixes. Each opening was played
twice with colors reversed:
\[
  \text{random opening}
  \;+\;
  \text{deterministic play}
  \;+\;
  \text{color reversal}.
\]
This protocol samples many positions without allowing move noise to dominate
the strength estimate. The pair controls much of opening and first-player
advantage. It is reproducible if the opening list is saved, without reusing one
live random seed everywhere.

\subsection{Strength checks measured what losses could not}

Every five hours, the loop played a 40-game paired-opening match against a
fixed reference. The run continued for the full budget even if several checks
were flat; the ladder was evidence, not a stopping rule.

Use validation loss to diagnose training distribution, overfitting, and
serialization. Use games to measure playing strength. These quantities can
move in opposite directions because:
\begin{itemize}
  \item policy targets are produced by changing searches;
  \item replay is non-stationary and correlated;
  \item stronger play may visit harder positions;
  \item cross-entropy rewards matching visit distributions, not winning by
        itself;
  \item value calibration and move ranking can improve differently.
\end{itemize}

\section{Bug prevention: make contracts executable}

\subsection{Tests worth writing before the first long run}
\begin{enumerate}
  \item \textbf{Reference rules.} Compare the optimized legal-action generator
        with a literal, independently written version on thousands of random
        states.
  \item \textbf{Transition round trip.} For every legal action in sampled
        states, make then unmake and compare complete snapshots.
  \item \textbf{Action round trip.} Require
        $\mathrm{decode}(\mathrm{encode}(a))=a$ for every legal action.
  \item \textbf{Terminal semantics.} Hand-construct each terminal cause and
        check both the winner and whether the player-to-move field changes.
  \item \textbf{Perspective.} Evaluate color-swapped or mover-swapped copies and
        verify the exact expected relation between relative and absolute value.
  \item \textbf{Symmetry.} Transform state, legal actions, policy, and outcome;
        invert the transform; require equality.
  \item \textbf{Hand-computed PUCT.} Build tiny trees for both players and check
        the chosen child, visits, $W$, and $Q$ after each backup.
  \item \textbf{Neural masking.} Assert exactly zero probability on illegal
        actions and unit probability mass over legal actions.
  \item \textbf{Serialization.} Save and reload a trained model; compare
        predictions and optimizer iteration/state, not just weights.
  \item \textbf{Actor identity.} Hash the model used to create every tranche and
        the model written after training. A completed update must change the
        expected hash.
\end{enumerate}

\subsection{Symptom-to-test debugging table}

\small
\begin{longtable}{L{0.25\textwidth} L{0.34\textwidth} L{0.33\textwidth}}
\toprule
\sffamily\bfseries Symptom &
\sffamily\bfseries Likely classes of cause &
\sffamily\bfseries Fastest discriminating tests \\
\midrule
\endfirsthead
\toprule
\sffamily\bfseries Symptom &
\sffamily\bfseries Likely classes of cause &
\sffamily\bfseries Fastest discriminating tests \\
\midrule
\endhead
One player collapses &
Value sign, terminal turn switch, directional encoding, asymmetric search
budget, missing symmetry. &
Hand PUCT tests for both players; color-swapped tactical pairs; sequential
versus batched search at a larger budget. \\
\addlinespace
Policy loss falls but play does not improve &
Targets too weak, illegal-mask mismatch, replay overuse, identical actor,
evaluation too noisy. &
Hash checkpoints; inspect root targets; play fixed-position search sweeps; run a
paired deterministic arena. \\
\addlinespace
Value is accurate on training but wrong in games &
Perspective conversion, leakage between train/validation, correlated split,
normalization-state mismatch. &
Split by complete game; evaluate swapped copies; compare training and inference
modes; reload before evaluation. \\
\addlinespace
Good tactics, terrible long games &
Replay lacks strategic phases, search too shallow, fast-turn policy drifts,
terminal reward sparse. &
Histogram target plies and phases; label sampled states with a stronger search;
increase budget on the same checkpoint. \\
\addlinespace
Two fresh evaluations are identical &
Reused seeds, reused opening file, deterministic actor launch, identical
checkpoint hashes. &
Log opening hashes, first moves, process seeds, and model hashes before the
match. \\
\addlinespace
Arena result is extreme or unstable &
Color/opening bias, time forfeits, unequal hardware, too few independent
openings, failures counted as losses. &
Reverse colors, report failures separately, compare clocks and hardware, and
use confidence intervals over opening pairs. \\
\addlinespace
Batched search disagrees with reference search &
Virtual-loss accounting, duplicate expansion, stale priors, wrong backup order,
state mutation. &
Use a deterministic mock network; compare per-edge visits after each batch;
hash state before and after search. \\
\addlinespace
Training suddenly stabilizes after restart &
Optimizer state lost, learning rate reset or frozen, replay path changed. &
Record optimizer iteration, learning rate, replay count, and a fixed prediction
vector before and after restart. \\
\bottomrule
\end{longtable}
\normalsize

\subsection{Log enough to reconstruct a failure}

For every training iteration, record:
\begin{itemize}
  \item actor checkpoint path and cryptographic hash;
  \item optimizer iteration and learning rate;
  \item new games, new positions, replay size, and training presentations;
  \item full/fast search budgets and fraction of stored turns;
  \item root-noise parameters and move-temperature schedule;
  \item outcomes by player, game length, target ply, policy entropy, and search
        value;
  \item policy loss, value loss, finite-value checks, and gradient or weight
        norms if available;
  \item arena opening IDs, colors, failures, clocks, hardware, and result.
\end{itemize}

Do not turn logging into a second framework. A CSV row and a short text summary
per iteration are sufficient if field names and units are stable.

\subsection{Use continuous diagnostics, not brittle tactical gates}

The final tactical suite contained 20 exact, balanced positions rather than six
hand-picked puzzles. Its value score was continuous:
\[
  \frac{1}{20}\sum_{i=1}^{20}
  v_{\mathrm{abs}}(s_i)\,z_{\mathrm{abs}}(s_i).
\]
Positive values mean the prediction tends to agree with the exact outcome;
magnitude reflects confidence. Color-swapped copies test a perspective
invariant; they are not additional independent evidence.

A tactical suite is a probe, not a complete strength measure. Keep the raw
per-position outputs so one easy family cannot hide a systematic failure.

\section{Adapting the method to another environment}
\label{sec:adapt}

Some conclusions below are mathematical; others exploit the game. Make this
separation before transferring the design.

\begin{center}
\small
\begin{tabularx}{\textwidth}{L{0.28\textwidth} Y}
\toprule
\sffamily\bfseries Breakthrough fact & \sffamily\bfseries Consequence here \\
\midrule
Two-player, alternating, zero-sum &
One scalar value and a parent-player sign are sufficient. \\
Deterministic, perfect information &
No chance nodes or belief states are required; paired openings control much of
the evaluation variance. \\
Players are equivalent after rotation and relabelling &
A two-plane mover-relative neural representation is exact. \\
Left and right are strategically symmetric &
Horizontal reflection is valid augmentation; action channels must reflect too. \\
At most three pawn moves leave a square &
Three policy channels per board square form a direct action encoding. \\
Games are short enough for terminal outcomes &
The final winner supplies an undiscounted target to every stored position. \\
\bottomrule
\end{tabularx}
\end{center}

If one row becomes false in a new environment, redesign the corresponding
component rather than preserving the Breakthrough implementation by habit.

\subsection{Decide the value object before writing MCTS}

For a two-player zero-sum game, a scalar value is enough. Choose either a fixed
global viewpoint or a mover-relative viewpoint and derive selection and backup
algebra on paper.

For a general-sum or multiplayer game, a scalar current-player value can discard
information. Store a value vector
\[
  \mathbf{v}(s) = (v_1(s),\ldots,v_k(s))
\]
and let the parent choose using the component belonging to its decision maker.
Chance nodes use expectations rather than maximization. In partially observable
games, the network input and tree state must represent an information state;
tests must explicitly prevent hidden-state leakage.

\subsection{Canonicalize only exact equivalences}

\begin{center}
\small
\begin{tabularx}{\textwidth}{L{0.21\textwidth} Y Y}
\toprule
\sffamily\bfseries Environment property &
\sffamily\bfseries Good adaptation &
\sffamily\bfseries Required test \\
\midrule
Players are interchangeable under a known transform &
Canonicalize to the mover or a fixed seat. &
Transform state, actions, and rewards; invert and recover the original sample. \\
Board has geometric symmetries &
Sample a group element during training or ensemble at inference if justified. &
Legal actions and transition results commute with every transform. \\
Rules contain asymmetric terrain, roles, or objectives &
Keep the asymmetry in the observation; do not augment it away. &
Construct a counterexample position where the proposed transform would change
the game. \\
Stochastic transitions &
Represent chance outcomes explicitly or train on sampled returns. &
Check probability mass, reproducibility under controlled RNG, and unbiased
backup estimates. \\
Very long horizon &
Use denser diagnostics, search-budget scheduling, and possibly reward/value
normalization. &
Test terminal rewards separately from any shaped reward and inspect phase
coverage in replay. \\
Continuous actions &
Replace full policy enumeration with a suitable proposal/search method. &
Test bounds, density normalization, sampling, and action-transform Jacobians if
needed. \\
\bottomrule
\end{tabularx}
\end{center}

\subsection{Match data generation to the environment's difficulty}

The successful recipe used rollout-PUCT pretraining because Breakthrough had a
cheap legal simulator and meaningful shallow tactics. Other environments may
need:
\begin{itemize}
  \item demonstrations to reach rare rewarding states;
  \item curriculum levels to control horizon or branching factor;
  \item an exact solver for small states;
  \item a heuristic or model-predictive controller as teacher;
  \item prioritized scenario generation for safety-critical or rare events.
\end{itemize}

Whichever source is chosen, measure the resulting policy, state coverage,
outcome balance, and effective sampling weights. One million rows says little
if they contain 10,000 copies of a few trajectories.

\subsection{Adapt evaluation to the source of variance}

Paired color-reversed openings are natural for deterministic, alternating,
two-player board games. In another environment, pair or stratify over the
largest nuisance variables:
\begin{itemize}
  \item random seeds and initial conditions in stochastic control;
  \item maps, scenarios, and traffic patterns in simulators;
  \item roles or teams in asymmetric games;
  \item demand regimes and time periods in operations problems.
\end{itemize}

Keep policy stochasticity out of evaluation unless stochastic execution is part
of the deployed agent. Report failures and timeouts separately from legitimate
losses.

\section{A staged implementation protocol}

The following order minimizes the cost of being wrong. Each gate should have a
written pass condition and a saved artifact or test result.

\begin{enumerate}
  \item \textbf{Environment gate.} Implement the smallest faithful ruleset, a
        literal reference transition/legal-action function, terminal rewards,
        and state cloning or make/unmake. Pass differential and round-trip
        tests.
  \item \textbf{Action and value gate.} Define action indexing, neural
        canonicalization, the public search value convention, and parent
        utility. Exhaust action round trips and player/symmetry metamorphisms.
  \item \textbf{Non-neural search gate.} Use uniform legal priors and a rollout
        or exact evaluator. Find immediate wins for every player, preserve the
        input state, and verify hand-computed $Q$ selection.
  \item \textbf{Network gate.} Train on solver-labelled positions.
        Require held-out learning, legal masking, save/load equality, and
        perspective consistency.
  \item \textbf{Teacher data gate.} Generate a reconstructable search
        corpus. Audit game uniqueness, seat balance, terminal legality, target
        mass, and validation split by whole episode.
  \item \textbf{One-cycle gate.} From one checkpoint: generate, add to replay,
        train with preserved optimizer state, save, reload, and generate again.
        Assert the actor hash changed and all metrics are finite.
  \item \textbf{Small-environment learning gate.} Run enough cycles to beat
        pretraining and fixed baselines in paired arenas. Diagnose regressions
        before scaling.
  \item \textbf{Native target gate.} Create a native network and new data.
        Transfer only measured process choices. Calibrate search budget before
        serious compute.
  \item \textbf{Serious-run gate.} Fix a compute budget, replay exposure, search
        schedule, exploration policy, strength-check protocol, recovery plan,
        and artifact paths before submission.
\end{enumerate}

\begin{keybox}
\boxtitle{Rule for expensive jobs}
Every expensive run should answer a question that a cheaper run cannot answer.
If the next 20 GPU-hours are merely checking whether action indices are correct,
the project has skipped a gate.
\end{keybox}

\section{Problems worth working through}

These exercises target the places where a plausible implementation most often
goes wrong. They are more useful than copying a complete training script.

\begin{enumerate}
  \item \textbf{PUCT by hand.} At a Player-2 parent, child $A$ has
        $Q=+0.30$ and $U=0.05$; child $B$ has $Q=-0.05$ and $U=0.20$.
        Compute both selection scores under the absolute convention. Then
        repeat for Player 1. Check: Player 2 chooses $B$; Player 1 chooses $A$.

  \item \textbf{Trace a terminal value.} Player 2 moves from $s_0$, Player 1
        moves from $s_1$, and Player 2 eventually wins. Write the absolute
        replay outcome and the two neural value targets. Check: the stored
        outcome is $-1$ for both rows; the mover-relative targets are $+1$ at
        $s_0$ and $-1$ at $s_1$.

  \item \textbf{Make symmetry executable.} Choose a legal capture near the
        left edge. Reflect the board, move tuple, legal mask, and a one-hot
        policy. Verify the transition commutation equation from Section 4 and
        exact recovery after two reflections. Make the test fail once by
        leaving the policy unchanged; inspect why ordinary shape checks miss
        the error.

  \item \textbf{Measure replay exposure.} An iteration adds 4,000 positions and
        draws 8,000 examples for training. Compute the reuse ratio. Then explain
        why this does not mean that every one of 25,000 replay entries was seen
        twice. Check: the ratio is 2 per \emph{fresh} position; sampling is over
        the whole window and individual counts vary.

  \item \textbf{Find validation leakage.} A 30-ply game contributes 30 adjacent
        positions, each optionally reflected. Compare a random position split
        with a whole-game split. List two ways the first can make validation
        optimistic, and state when the validation set must be chosen.

  \item \textbf{Design a fair 40-game arena.} Specify how many independently
        generated opening prefixes are needed, how colors are assigned, what
        randomness remains after the opening, which model hashes are logged,
        and how crashes are reported. Check: 20 openings, each played twice;
        deterministic search after the opening; failures separate from game
        results.

  \item \textbf{Discriminate causes.} A new batched actor shows a severe color
        imbalance. Propose the smallest set of fixed-checkpoint runs that can
        distinguish a value-sign bug, a batching bug, and insufficient search.
        Each comparison should change only one relevant feature.
\end{enumerate}

\begin{plainbox}
\boxtitle{Implementation lab}
Before launching 8x8 self-play, make all seven exercises concrete on 5x5:
print the hand-computed node statistics, save one audited replay row, and keep
the failing version of each metamorphic test long enough to see that it detects
the intended defect. A test whose failure you have never observed may be
checking less than its name suggests.
\end{plainbox}

\section{The successful Breakthrough recipe}

These settings describe the successful regime. The relationships are more
portable than the exact numbers.

\small
\begin{longtable}{L{0.20\textwidth} L{0.25\textwidth} L{0.42\textwidth}}
\toprule
\sffamily\bfseries Component &
\sffamily\bfseries Breakthrough setting &
\sffamily\bfseries Portable principle \\
\midrule
\endfirsthead
\toprule
\sffamily\bfseries Component &
\sffamily\bfseries Breakthrough setting &
\sffamily\bfseries Portable principle \\
\midrule
\endhead
Neural input &
Two mover-relative planes. &
Canonicalize only under an exact, reversible transform for state, action, and
reward. \\
Network &
48 filters, 3 residual blocks, no BatchNorm. &
Use the smallest model that passes supervised and save/load gates. Enlarge only
after data and search are credible. \\
Pretraining &
10,000 rollout-PUCT games; 421,561 stored positions. &
Begin neural self-play with plausible search targets; measure a data-size curve
rather than guessing corpus sufficiency. \\
Replay &
Recent FIFO, about 25,000 neural positions. &
Balance on-policy freshness against trajectory and skill diversity. \\
Optimization &
Batch 128, reuse ratio about 2, Adam retained, learning rate
$2.5\times10^{-4}$. &
Track presentations per fresh position and preserve optimizer state. \\
Augmentation &
Random left-right reflection after sampling. &
Store real samples once; transform state and policy together during training. \\
Self-play exploration &
Root Dirichlet noise and visit sampling for the first 12 plies. &
Explore early enough to diversify trajectories, then sharpen to finish games. \\
Parallelism &
32 games advanced together; neural leaves batched. &
Batch inference without changing search semantics; compare with a slow
reference implementation. \\
Search economy &
25\% full-search stored targets; fast intervening turns. &
Spend computation on the positions that become labels; audit phase coverage. \\
Search schedule &
128/16, then 192/24, then 256/32 full/fast simulations. &
Calibrate the minimum useful budget on a fixed checkpoint and revisit as the
model improves. \\
Evaluation &
40 games every 5 hours; distinct six-ply openings, colors reversed, no search
noise. &
Use deterministic agents over diverse paired initial conditions. Keep training
independent of a noisy acceptance gate. \\
\bottomrule
\end{longtable}
\normalsize

\subsection{Anti-patterns to avoid}
\begin{itemize}
  \item Do not change value perspective in several modules.
  \item Do not apply a symmetry to observations without applying its exact
        action permutation to policy targets.
  \item Do not infer corpus size from augmented presentations; count unique
        stored states and complete games.
  \item Do not call a random position split validation when adjacent states
        from the same game appear in both sets.
  \item Do not treat a near-zero learning rate as a cure for bad labels.
  \item Do not conclude that batching is wrong before repeating the test at a
        useful search budget.
  \item Do not let failed games, timeouts, or identical checkpoints enter an
        arena score silently.
  \item Do not reuse one fixed seed throughout production and mistake
        reproducibility for diversity.
  \item Do not scale to the target environment because a reduced architecture
        happens to accept padded inputs.
  \item Do not add gates, rollback rules, or distributed services faster than
        you can test their identities and state transitions.
\end{itemize}

\section{Readiness checklist}

\begin{center}
\small
\begin{tabularx}{\textwidth}{L{0.28\textwidth} Y}
\toprule
\sffamily\bfseries Before launching &
\sffamily\bfseries Evidence required \\
\midrule
Rules and state &
Differential legal-action tests; complete make/unmake snapshots; all terminal
causes covered. \\
Action representation &
Exhaustive legal encode/decode round trips; legal-mask agreement. \\
Perspective &
One written convention; both-player hand PUCT tests; swapped-position
metamorphic tests. \\
Symmetry &
State/action/policy/reward transform tests and inverse recovery. \\
Network &
Exact-label supervised learning; illegal probability zero; save/load prediction
and optimizer equality. \\
Data &
Unique game count, position count, outcome/seat balance, whole-game split,
effective weights, and augmentation policy recorded. \\
Search &
Immediate tactics for each player; state preserved; fixed-checkpoint budget
sweep; reference versus batched agreement. \\
Training cycle &
Actor hash, replay growth, finite losses, optimizer continuity, and changed
post-update hash. \\
Evaluation &
Distinct initial conditions, pairing/stratification, equal resources, failure
count, model hashes, and uncertainty interval. \\
Recovery &
Atomic checkpoint, replay location, iteration log, resume test, and known last
complete unit of work. \\
\bottomrule
\end{tabularx}
\end{center}

\section{What we would do next}

Declaring success does not end the experiment. The remaining 50-hour run maps
the learning curve: whether improvement continues, saturates, oscillates, or
regresses as the replay distribution moves.

The most valuable follow-up is not immediately a larger network. It is a small
ablation matrix around the successful recipe:
\begin{itemize}
  \item hold the checkpoint fixed and compare search budgets;
  \item train matched short continuations with and without playout-cap
        randomization;
  \item compare two replay windows at equal fresh-position reuse;
  \item compare retained and reset optimizer state;
  \item repeat arenas across hardware with simulation-count rather than
        wall-clock limits.
\end{itemize}

These experiments turn causal hypotheses into measured claims. They also give
students a more useful result than a claim that the final settings worked: a
picture of which parts are robust and which depend on the environment and
compute regime.

\section{Conclusion}

The successful project was not rescued by a mysterious neural-network trick.
It became learnable when the entire feedback loop agreed with itself:
\begin{enumerate}
  \item the environment produced correct states and rewards;
  \item action, symmetry, and perspective transformations were explicit;
  \item search used the value convention intended by training;
  \item the starting data were large and competent enough to matter;
  \item replay and optimization exposed fresh targets at a controlled rate;
  \item parallelism accelerated inference without redefining search;
  \item self-play explored, while evaluation controlled its sources of noise;
  \item strength was measured by paired games and checkpoint identity was
        verified.
\end{enumerate}

\begin{keybox}
\boxtitle{The transferable lesson}
An AlphaZero implementation is a circular scientific instrument. Every arrow in
the circle needs a unit, a viewpoint, an identity, and a test. Once those
contracts are trustworthy, the sophisticated mathematics can do its job. Before
that, more compute simply makes the wrong experiment run faster.
\end{keybox}

\appendix
\section{Failure timeline and evidence}
\begin{itemize}
  \item A prior continuous 8x8 run was invalidated by symmetry-state loss across
        learner processes and by an omitted inference ensemble.
  \item A large-data diagnosis found rare teacher examples with roughly 71 times
        the per-position weight of neural examples and effective batches near
        18 despite a nominal size of 256.
  \item A supposedly meaningful arena was shown invalid when both checkpoint
        hashes were identical; another was invalidated by time forfeits.
  \item The clean-room project established 5x5 correctness and learning before
        creating native 8x8 weights and data.
  \item The 1,000-game pretrained model lost 24--76 to the 10,000-game model,
        showing that corpus scale mattered in this regime.
  \item Phase-6 iteration 5 beat the raw pretrained initializer 52--8 but lost
        2--58 to alpha-beta on RTX 2080-class hardware.
  \item Serious iteration 15 beat phase-6 iteration 5 by 39--1.
  \item Serious iteration 26 beat iteration 15 by 36--4.
  \item Serious iteration 26 beat depth-4 alpha-beta 49--1 on an L40S, using 25
        distinct six-ply openings with colors reversed.
\end{itemize}

\section{Interpretive limits}

The strongest evidence in this document is comparative, not absolute. The
iteration ladder uses controlled paired openings. The final alpha-beta result
is decisive for its recorded protocol, but its magnitude cannot be compared
directly with the earlier phase-6 match because the hardware differed under a
wall-clock time control.

The project did not run a factorial ablation over all serious-loop ingredients.
Claims about why the bundle worked are therefore engineering judgments informed
by failed runs, unit tests, diagnostics, and measured matches. They should guide
the next experiment, not substitute for it.

Finally, Breakthrough is deterministic, perfect-information, alternating,
two-player, and zero-sum. The value algebra, evaluation pairing, and canonical
player transform exploit those facts. Section~\ref{sec:adapt} describes the
changes required when an environment does not share them.

\section{Further reading}

\begin{itemize}
  \item Dominik Klein, \emph{Neural Networks for Chess}, especially the MCTS
        chapter, the AlphaZero overview, and the executable HexapawnZero case
        study. \href{https://arxiv.org/abs/2209.01506}{arXiv:2209.01506}.
  \item David Silver et al., \emph{Mastering Chess and Shogi by Self-Play with
        a General Reinforcement Learning Algorithm},
        \href{https://arxiv.org/abs/1712.01815}{arXiv:1712.01815}. This is the
        primary source for the AlphaZero algorithm and its chess/shogi
        instantiation.
  \item David J. Wu, \emph{Accelerating Self-Play Learning in Go},
        \href{https://arxiv.org/abs/1902.10565}{arXiv:1902.10565}. KataGo is a
        useful example of treating self-play efficiency, target quality, and
        evaluation as engineering questions rather than historical dogma.
  \item Richard S. Sutton and Andrew G. Barto, \emph{Reinforcement Learning:
        An Introduction}, second edition, for the broader language of value
        functions, bootstrapping, exploration, and off-policy data.
\end{itemize}

Klein's Hexapawn chapter is particularly useful before this guide: it shows the
complete loop in a game small enough to reason about. This post-mortem takes
the next step---what must be added when the implementation is large enough that
silent interface errors, non-stationary data, parallel inference, and noisy
strength estimates become decisive.

\end{document}
